\section{Introduction}
\label{sec:intro}

Tabular classification remains dominated by gradient-boosted decision trees
\citep{grinsztajn2022tree,mcelfresh2023neural}, yet no single algorithm wins
everywhere: the best model varies from dataset to dataset, and choosing it by
cross-validation is itself noisy. Ensembling a diverse library of models
sidesteps the choice. Greedy ensemble selection \citep{caruana2004ensemble} and
multi-layer stacking \citep{wolpert1992stacked} are the workhorses behind
leading AutoML systems such as auto-sklearn \citep{feurer2015efficient} and
AutoGluon \citep{erickson2020autogluon}, but inside those systems they are
entangled with hyperparameter search, time budgeting and dozens of engineering
choices, which makes it hard to tell which ingredients matter.

NEXUS is a deliberately small, fully specified ensemble pipeline: a fixed
library of eight algorithms with a~priori hyperparameters, bagged cross-fitting,
an optional second stacking layer, greedy selection of convex ensemble weights,
and a single calibration temperature. It contains no search and no time
budget, so its behaviour is deterministic and every component can be ablated.

Our contributions are empirical and methodological rather than a new learning
algorithm:
\begin{enumerate}
\item \textbf{A held-out design protocol.} Every design decision for the
final method was made on 16 development datasets disjoint from the evaluation
suite, and frozen (and committed to version control) before the method was run
on the test suite. \Cref{sec:dev} reports all eleven candidates we tried,
including the ones that did not help.
\item \textbf{A large, statistically careful evaluation} on \NTasks{} tasks of
OpenML-CC18 using the official 10-fold splits, against eight single models, a
tuned gradient-boosting baseline, three ensembling strategies on identical
base predictions, and AutoGluon given the same per-fold time budget.
\item \textbf{Component-level evidence}: ablations of every model family and of
calibration, the effect of ensemble size, and a measurement of a systematic
bias in out-of-fold model selection: bagging improves high-variance learners
(e.g.\ $k$-NN, LightGBM) far more than low-variance ones, so out-of-fold
estimates misjudge their bagged test-time quality.
\item \textbf{A reproducible open-source package}: a scikit-learn compatible
estimator, the benchmark harness, raw-prediction caches, and a single script
that regenerates every number, table and figure in this paper.
\end{enumerate}
